% !TEX root = ../thesis.tex

\chapter{CFD 高保真验证模板}
\label{chap:wind_cfd}

本章给出本研究中 CFD 高保真验证模板的搭建细节，
对应开题报告第四项研究任务（T4）。
中期阶段的研究场景限定为恒定自由来流条件下的稳态前飞工况优化（详见第~\ref{sec:scope}~节），
因此本章不涉及湍流风场建模与风场—优化耦合（这部分作为毕业论文阶段的工作）。
本章首先给出基于 Star-CCM+ 的 CFD 高保真验证模板的物理模型与计算域设置；
随后介绍网格划分策略、边界层处理与湍流模型选择；
最后说明验证策略从 CFD 主导调整为 QBlade 回验为主、CFD 辅助高保真的依据。

\section{基于 Star-CCM+ 的 CFD 验证模板}

\subsection{建模动机}

为加强优化结果的物理可解释性，本研究在 LLFVW 同源回验之外、
额外搭建 Star-CCM+ 商用 CFD 平台作为辅助高保真验证手段，
对优化得到的代表性外形开展全工况扫描，
并在"基准桨（APC 10$\times$7）与代表性优化桨"上完成对照渲染。
相比于 LLFVW 同源回验，
CFD 能够显式解析螺旋桨表面边界层、桨尖涡卷起、近场尾迹结构等三维黏性效应，
对于代理模型在低保真度（LLFVW）数据上得到的优化结果，
CFD 给出独立的高保真度交叉校核，
从而提高最终设计结果的物理可信度。

\subsection{物理模型与计算域}

CFD 验证模板的物理模型搭建与流体域划分遵循以下原则：
理论上，计算区域越大，计算域边界对数值模拟结果的影响就越小；
然而，过大的计算域会导致计算资源的浪费，
而过小的计算域则会导致螺旋桨的气动特性难以被精确捕捉。
因此，在选择计算域时需要遵循两条原则：

\begin{enumerate}
  \item 尽量减小计算域对流场的影响，
        以便使尾部的湍流区得以充分发展；
  \item 计算域的阻塞比 $\phi$ 应小于 $5\%$。
\end{enumerate}

其中，阻塞比 $\phi$ 定义为
\begin{equation}
\phi = \frac{A_1}{A},
\end{equation}
其中 $A_1$ 为螺旋桨圆盘的正投影面积，$A$ 为数值模拟入口处的面积。

根据上述原则，计算域的物理定义与尺寸大小为：
入口区与出口区直径为 $10D$（$D$ 为螺旋桨直径），
旋转域直径为 $1.2D$、宽度为 $0.4D$；
计算域总宽度为 $25D$；
共产生约 700 万网格，
阻塞比为 $2.5\%$，
满足 $\phi < 5\%$ 的工程要求。

\subsection{网格策略与边界层}

边界层网格基于 $y^+$ 函数的近似，
保证 $y^+$ 区间满足 SST $k$-$\omega$ 湍流模型的适用范围。
最终生成的 $y^+$ 直方图显示其取值落在 $0 \sim 1$ 之间，
完全满足壁面分辨率要求。

\subsection{湍流模型与精度目标}

CFD 验证模板的关键参数汇总如表~\ref{tab:cfd_template}~所示。
根据文献\cite{Goetten2019}（亚音速气动 RANS CFD 仿真指南综述）给出的参数进行设置，
所得功率系数 $C_P$、推力系数 $C_T$ 与推进效率 $\eta$ 的相对误差均控制在 $5\%$ 以内，
与文献结果一致，
满足优化结果验证及可视化需求。

\begin{table}[htbp]
\centering
\caption{Star-CCM+ CFD 验证模板参数}
\label{tab:cfd_template}
\begin{tabular}{ll}
\toprule
\textbf{参数} & \textbf{值} \\
\midrule
计算域 & 入口 / 出口 $10D$，总宽 $25D$ \\
旋转域 & 直径 $1.2D$，宽 $0.4D$ \\
网格量 & $\sim 700$ 万 \\
阻塞比 $\phi$ & $2.5\%$（$< 5\%$） \\
湍流模型 & SST $k$-$\omega$ \\
$y^+$ 范围 & $0 \sim 1$（满足壁面分辨率要求） \\
精度目标 & $C_P$、$C_T$、$\eta$ 相对误差 $\leq 5\%$ \\
\bottomrule
\end{tabular}
\end{table}

\subsection{基线桨对照渲染}

为验证模板的工程可用性，
已在 APC 10$\times$7 基线桨上完成 1$\sim$2 个代表性工况的 CFD 仿真，
并将结果与 LLFVW 仿真在相同工况下的输出做对照：

\begin{itemize}
  \item \textbf{性能参数一致性：}
        $C_T$、$C_P$、$\eta$ 相对误差均 $\leq 5\%$；
  \item \textbf{尾迹结构：}
        CFD $Q$-criterion 等值面识别得到的桨尖螺旋涡与 LLFVW 自由涡环位置基本一致，
        两种方法在桨尖涡卷起阶段的几何形态吻合；
  \item \textbf{压力分布：}
        CFD 给出的桨叶表面压力分布与 LLFVW 升力分布在径向积分上一致，
        在局部分布上 CFD 能解析根部分离区等 LLFVW 无法捕捉的细节。
\end{itemize}

对照结果确认 CFD 模板可作为优化最终解的独立高保真验证手段。

\section{验证策略：QBlade 回验为主、CFD 辅助}

\subsection{开题原计划：CFD 主导终验}

开题报告中规划的验证路径为：
LLFVW 优化得到代表性外形 $\rightarrow$
CFD 高保真仿真终验 $\rightarrow$
$Q$-criterion 等可视化分析。
该路径将 CFD 作为可信度的最终裁决标准。

\subsection{实际调整：QBlade 回验优先，CFD 辅助}

在实际推进过程中发现，
单次 CFD 仿真耗时数天且对几何细节高度敏感，
若每轮优化结果均需 CFD 终验，
则优化—验证闭环周期过长，
难以支撑主动补点循环的高频迭代。
因此，验证策略调整为：

\begin{enumerate}
  \item \textbf{QBlade LLFVW 回验为主}：
        优化得到的代表性外形首先在 QBlade LLFVW 框架内进行全工况扫描验证
        （同源闭环，与训练数据同分布，计算成本低），
        筛选出通过同源验证的候选外形进入 CFD 终验阶段；
  \item \textbf{Star-CCM+ CFD 辅助高保真验证}：
        模板已就绪并通过基线桨对照（误差 $\leq 5\%$），
        在 QBlade 回验通过后对 1$\sim$2 个最终优化解开展代表性工况的 CFD 终验，
        作为高保真结果支撑，
        并通过 $Q$-criterion、速度切面、压力切面等可视化手段
        分析优化前后的尾迹涡结构差异。
\end{enumerate}

该调整在不损失高保真验证强度的前提下，
显著提升了优化迭代的吞吐率，
为主动补点闭环的高频运行提供了支撑。

\section{本章小结}

本章给出了基于 Star-CCM+ 的 CFD 验证模板的完整搭建细节
（$\sim 700$ 万网格、$y^+ \in [0, 1]$、SST $k$-$\omega$、误差 $\leq 5\%$），
并在 APC 10$\times$7 基线桨上完成了模板验证，
确认其工程可用性。
基于实际研究中遇到的 CFD 计算周期约束，
验证策略由 CFD 主导调整为 QBlade 回验优先 + CFD 辅助高保真，
在不损失验证强度的前提下显著提升了优化—验证闭环的迭代效率。
中期阶段，
本课题的研究场景限定为恒定自由来流条件下的稳态前飞工况优化；
湍流风场建模与风场—优化耦合作为毕业论文阶段的扩展工作，
不在本中期论文范围之内。
